\begin{enumerate}
	\item Les Pays-Bas ont obtenu 27 médailles d'or.
	\item $63-17-28=18$.

L'Australie a obtenu 18 médailles d'or.
	\item La Grande-Bretagne a obtenu 31 médailles de bronze, sur un total de 124.

		\begin{tabular}{c|c|c}
		&en nombre  &en \% \\\hline
médailles de bronze		& 31 &$x$  \\\hline
toutes les médailles&124&100\\
	\end{tabular}


	$x=\dfrac{31\times100}{124}=25$: le pourcentage de médailles de bronze est 25~\%, donc l'affirmation est vraie.
	\item 	\begin{enumerate}
				\item Classons ces nombres par ordre croissant :

				$56 ~<~ 63 ~<~ 71 ~<~ 75 ~<~ \textcolor{red}{82} ~<~ 89 ~<~ 105 ~<~ 124 ~<~ 220$

				La médiane est la valeur \og centrale \fg{}, c'est-à-dire la 5\up{ème} : 82 médailles.
				\item Au moins la moitié de ces neuf pays ont eu 82 médailles ou moins de 82 médailles et au moins la moitié ont eu 82 médailles ou plus de 82 médailles.
			\end{enumerate}
	\item 20 médailles représentent nos 100\% de départ.

	\(\begin{array}{ccc}
	20 \text{ médailles}& \to & 100\%\\
	26 \text{ médailles}& \to & ?
	\end{array}\)

	$?=\dfrac{100\times 26 }{20}=130$.

	$130-100=30$

	Le nombre de médailles obtenues par le Brésil a augmenté de 30\%.
\end{enumerate}

\needspace{10\baselineskip}
% saut de page si moins de 10 lignes disponibles
